\documentclass[conference]{IEEEtran}
% IEEEtran sets \rmdefault to ptm (Times), which the IEEE requires. Load real
% Times-metric fonts explicitly so every engine (pdflatex on Overleaf, tectonic
% locally) lays the text out identically instead of falling back to Latin Modern.
\IEEEoverridecommandlockouts
\usepackage{cite}
\usepackage{amsmath,amssymb,amsfonts}
% newtx must come AFTER amssymb (both define \Bbbk)
\usepackage{newtxtext,newtxmath}
\usepackage{graphicx}
\usepackage{booktabs}
\usepackage{textcomp}
\usepackage{xcolor}
\usepackage{url}
\def\BibTeX{{\rm B\kern-.05em{\sc i\kern-.025em b}\kern-.08em
    T\kern-.1667em\lower.7ex\hbox{E}\kern-.125emX}}

\begin{document}

\title{Zero-Shot Coverage of Unseen Substance\\
Classes for SUD Treatment-Dropout\\
Prediction on 6.7 Million Treatment Episodes}

% UGHS requirement: the affiliation must indicate whether the FIRST author is a high
% school or an undergraduate student. Per author decision the explicit designation line
% was removed; the school name in the affiliation carries it implicitly. The 5-page
% high-school limit still governs this manuscript.
% ACTION FOR AUTHORS: replace the email placeholder below before submission.
\author{\IEEEauthorblockN{Raymond Cui}
\IEEEauthorblockA{\textit{White Station High School}, Memphis, Tennessee, USA\\
\texttt{<author email>}}}

\maketitle

\begin{abstract}
Predicting substance-use-disorder (SUD) treatment dropout is well studied, yet models are
typically trained one substance at a time, leaving small or emerging classes uncovered. We ask
whether a single model trained only on the six highest-volume substances can predict usefully for
classes it has never seen: zero-shot in the sense that no training episode belongs to the
target class. Using 6.72 million episodes from the U.S. Treatment Episode Data Set (TEDS-D,
2015--2022), we train a pooled gradient-boosted tree on 15 intake-time features and evaluate it
on twelve held-out small classes. It achieves mean AUROC 0.777 (range 0.745--0.806) on unseen
classes, nearly matching performance on trained substances, while pooling yields no accuracy gain
over substance-specific training. On those unseen classes it stays calibrated (mean decile gap
0.015) and clinically useful (positive net benefit to threshold 0.80), loses more than the trained
substances across the COVID-19 discontinuity (0.051 versus 0.017), and audits clean by subgroup
except at age 65+, where predicted risk exceeds observed by 0.140. TabPFN-3 is more
sample-efficient at matched budgets, yet at each model's maximum training size both converge to a
shared $\sim$0.78 AUROC ceiling, locating the limit in the intake features rather than the
learner. One deployable model therefore covers the substance long tail without per-class
retraining.
\end{abstract}

\begin{IEEEkeywords}
substance use disorder, treatment retention, machine learning, gradient boosting,
zero-shot generalization, tabular foundation models, TabPFN, health informatics
\end{IEEEkeywords}

\section{Introduction}
Retaining patients in substance-use-disorder (SUD) treatment is one of the strongest modifiable
predictors of recovery, and dropout is common: in the national data used here, 39.8\% of resolved
episodes end in the patient leaving against professional advice. If the episodes at highest risk of dropout
could be flagged at intake, programs could target limited follow-up and retention resources more
efficiently, an application in which machine learning has shown broad promise across psychiatry
\cite{chekroud_ml_psych}.

Machine-learning models that predict SUD treatment completion or dropout from administrative
intake data are well established, including on TEDS-D
\cite{ml_teds_disparities,ml_jamia,attrition2025,sud_ml_review}, but essentially all of this
work trains and evaluates \emph{one substance at a time} or on a single mixed cohort, and never
asks whether a model generalizes to a substance class absent from training. National treatment
systems see a heavy-tailed substance distribution: a handful of high-volume substances dominate,
while dozens of smaller classes (and periodically new ones such as novel synthetic opioids)
individually have far too few episodes to train a reliable substance-specific model, so a
per-substance paradigm leaves them uncovered.

We therefore invert the usual question, asking not ``how accurately can we predict dropout for
substance $X$?'' but \textbf{can one model, trained only on the data-rich substances, produce
useful dropout predictions for the substance classes it has never seen?}

Our contributions are:
\begin{enumerate}
\item \textbf{Single-model coverage of unseen substance classes (primary).} A pooled model
trained on six substances predicts dropout for twelve held-out small classes at mean AUROC 0.777
with no substance-specific training, and we test calibration, clinical utility, temporal
robustness and subgroup fairness \emph{on those unseen classes}.
\item \textbf{Evidence that coverage, not accuracy, is the reason to pool.} Pooled training does
not improve on substance-specific training, so the pooled model earns its place by covering
classes a per-substance approach cannot reach rather than by scoring higher on the ones it can.
\item \textbf{A two-protocol tabular-foundation-model benchmark.} We benchmark TabPFN-3
\cite{tabpfn_nature} against gradient boosting at matched training size and at each model's
maximum training size, separating sample efficiency from attainable accuracy and exposing a
task-imposed ceiling.
\end{enumerate}

\section{Related Work}
\textbf{Dropout prediction on SUD administrative data.} On TEDS, Baird et al.\
\cite{ml_teds_disparities} used a virtual-twins random forest to study disparities in
completion, Rabinowitz et al.\ \cite{attrition2025} predicted attrition in one cohort
($N \approx 30$k) at AUROC 0.63--0.67, and Nasir et al.\ \cite{ml_jamia} reported
AUROC $\approx$0.89 for completion by joining discharge records with county-crime data. A
systematic review \cite{sud_ml_review} flags the shared gap: validation within a single substance or cohort, with out-of-distribution
generalization rarely tested.

\textbf{Transfer across substances and tabular foundation models.} The closest precedent is
Bailey and DeFulio \cite{bailey2022}, who transferred a convolutional network's weights from
a heroin dataset to a smaller opioid dataset, one source$\rightarrow$target pair between two
opioids, framed as a transfer method for an accuracy gain. Our design differs in scope and
conclusion: twelve held-out classes spanning sedatives, stimulants and hallucinogens, evaluated
with no weight adaptation, and the \emph{opposite} headline, that pooling does not raise
accuracy, so the value is coverage. Separately, TabPFN \cite{tabpfn_nature} is a transformer pretrained to
perform Bayesian inference over tabular datasets in one forward pass, strong in the low-data
regime; it has been applied to clinical prediction but, to our knowledge, not to SUD treatment
outcomes and not benchmarked against gradient boosting on a national administrative dataset.

\section{Data and Methods}

\subsection{Dataset}
We use the SAMHSA Treatment Episode Data Set, Discharges (TEDS-D), a public, de-identified,
national census of discharges from publicly funded SUD treatment facilities, pooling seven annual
files (2015--2020, 2022). The 2021 file was excluded because it ships text-labeled categories rather
than the numeric codes used in every other year, and a hand-built crosswalk would carry
miscoding risk; the retained years span pre-COVID, COVID-onset and post-COVID periods. After restricting to episodes with a resolvable outcome and a mappable primary
substance, the analytic sample is \textbf{6{,}717{,}226 episodes} with an overall dropout rate of
0.398.

\subsection{Task definition}
The prediction target is binary: \emph{dropout} (discharge reason ``left against professional
advice'') versus \emph{treatment completed}. Episodes with other discharge reasons (transfer,
incarceration, death, facility termination) are excluded, leaving a clean contrast.

Predictions use the \textbf{15 intake-time features} in Table~\ref{tab:features}, all recorded at
admission. The primary substance (SUB1) is only the grouping variable defining the coverage split
and is \textbf{never a feature}. To prevent leakage we exclude the discharge reason, all discharge (``\_D'')
fields, length of stay and the 18 substance flags that re-encode SUB1. In a single-year sensitivity check (2020 held out), expanding the panel to 27 pre-admission
features improved pooled AUROC by only $+0.005$ (0.820 to 0.825), within confidence-interval
width, so the 15-feature set was retained.

\begin{table}[t]
\caption{The 15 intake-time predictors. All are recorded at admission, so a prediction is
available before treatment begins. SUB1 (primary substance) serves only as the grouping
variable defining training strategies and is never a model input.}
\label{tab:features}
\centering
\scriptsize
\setlength{\tabcolsep}{3pt}
\begin{tabular}{@{}llp{0.46\columnwidth}@{}}
\toprule
TEDS field & Domain & Definition \\
\midrule
SUB2, SUB3 & Substance & Secondary and tertiary substances at admission \\
ROUTE1 & Substance & Route of administration (oral, smoking, inhalation, injection) \\
FREQ1 & Substance & Frequency of use of the primary substance \\
FRSTUSE1 & Substance & Age at first use of the primary substance \\
SERVICES & Treatment & Service setting (detox, residential, ambulatory) \\
DSMCRIT & Clinical & DSM diagnostic category \\
NOPRIOR & Treatment & Number of prior treatment episodes \\
PSOURCE & Referral & Referral source (self, provider, criminal justice) \\
AGE, GENDER & Demographic & Age band (12--14 to 65+) and reported sex \\
RACE & Demographic & Reported race/ethnicity \\
EDUC, EMPLOY & Demographic & Education completed and employment status \\
STFIPS & Geography & Reporting state (FIPS code) \\
\bottomrule
\end{tabular}
\end{table}

\begin{table*}[t]
\caption{Cohort characteristics of the analytic sample ($N=6{,}717{,}226$ discharge
episodes, 2015--2022, 2021 excluded). \%~is the column percentage of the full cohort and
\emph{Drop.} the within-group dropout rate; the cohort-wide rate is 0.398. The six
highest-volume substances (alcohol through cocaine) form the big-6 training pool and the
remaining twelve are the held-out coverage targets. Race is missing for 2.1\% of episodes
and sex for 0.04\%; age band is complete.}
\label{tab:cohort}
\centering
\scriptsize
\setlength{\tabcolsep}{2.5pt}
\begin{tabular}{@{}lrrr@{\hspace{9pt}}lrrr@{\hspace{9pt}}lrrr@{}}
\toprule
Group & $n$ & \% & Drop. & Group & $n$ & \% & Drop. & Group & $n$ & \% & Drop. \\
\midrule
\multicolumn{4}{@{}l@{}}{\textit{Primary substance}} & \multicolumn{4}{@{}l@{}}{\textit{Age band at admission}} & \multicolumn{4}{@{}l@{}}{\textit{Race/ethnicity}} \\
Alcohol & 2,491,184 & 37.1 & 0.295 & 12-14 & 35,002 & 0.5 & 0.426 & Alaska Native & 16,029 & 0.2 & 0.288 \\
Heroin & 1,560,370 & 23.2 & 0.474 & 15-17 & 182,273 & 2.7 & 0.432 & American Indian & 162,380 & 2.4 & 0.319 \\
Cannabis & 840,885 & 12.5 & 0.456 & 18-20 & 215,916 & 3.2 & 0.463 & Asian/Pacific Islander & 2,487 & 0.0 & 0.349 \\
Methamphetamine & 776,204 & 11.6 & 0.444 & 21-24 & 588,001 & 8.8 & 0.434 & Black/African American & 1,190,339 & 17.7 & 0.425 \\
Other opioids & 456,633 & 6.8 & 0.473 & 25-29 & 1,119,941 & 16.7 & 0.427 & White & 4,374,536 & 65.1 & 0.390 \\
Cocaine & 375,967 & 5.6 & 0.442 & 30-34 & 1,108,153 & 16.5 & 0.417 & Asian & 47,345 & 0.7 & 0.385 \\
Benzodiazepines & 76,380 & 1.1 & 0.334 & 35-39 & 903,291 & 13.4 & 0.407 & Other single race & 602,916 & 9.0 & 0.448 \\
Other amphetamines & 39,015 & 0.6 & 0.421 & 40-44 & 662,231 & 9.9 & 0.392 & Two or more races & 147,044 & 2.2 & 0.459 \\
OTC & 37,578 & 0.6 & 0.531 & 45-49 & 606,764 & 9.0 & 0.370 & Native Hawaiian/Pacific Islander & 32,251 & 0.5 & 0.394 \\
PCP & 18,098 & 0.3 & 0.401 & 50-54 & 577,785 & 8.6 & 0.347 & Missing & 141,899 & 2.1 & 0.221 \\
Non rx methadone & 9,867 & 0.1 & 0.485 & 55-64 & 625,069 & 9.3 & 0.323 & \multicolumn{4}{@{}l@{}}{\textit{Discharge year}} \\
Other hallucinogens & 8,096 & 0.1 & 0.398 & 65+ & 92,800 & 1.4 & 0.290 & 2015 & 986,614 & 14.7 & 0.383 \\
Barbiturates & 7,695 & 0.1 & 0.430 & \multicolumn{4}{@{}l@{}}{\textit{Sex}} & 2016 & 1,015,960 & 15.1 & 0.388 \\
Other stimulants & 7,624 & 0.1 & 0.402 & Male & 4,492,019 & 66.9 & 0.378 & 2017 & 1,080,620 & 16.1 & 0.383 \\
Other tranquilizers & 5,187 & 0.1 & 0.170 & Female & 2,222,621 & 33.1 & 0.437 & 2018 & 1,034,079 & 15.4 & 0.386 \\
Other sedatives hypnotics & 3,214 & 0.0 & 0.399 & \multicolumn{4}{@{}l@{}}{\textit{Substance group}} & 2019 & 1,042,783 & 15.5 & 0.406 \\
Inhalants & 2,607 & 0.0 & 0.427 & Big-6 (trainable) & 6,501,243 & 96.8 & 0.397 & 2020 & 782,004 & 11.6 & 0.426 \\
Other sedatives & 622 & 0.0 & 0.476 & Small (held out) & 215,983 & 3.2 & 0.404 & 2022 & 775,166 & 11.5 & 0.425 \\
\bottomrule
\end{tabular}
\end{table*}

\subsection{Cohort and substance split}
Table~\ref{tab:cohort} summarizes the cohort across all eighteen primary-substance classes and
the demographic strata. The six ``big-6'' substances have enough episodes to train a
substance-specific model; the remaining twelve ``small'' classes are the coverage targets and
together are only 3.2\% of episodes. The smallest big-6 substance (cocaine, 376k) is
4.9$\times$ larger than the largest small class (benzodiazepines, 76k). Base rates vary far more
across substances (0.170 to 0.531) than the cohort-wide 0.398 suggests, so one model must span
markedly different populations.


\subsection{Models and evaluation}
The primary model is a histogram-based gradient-boosted-tree classifier (GBT)
\cite{friedman_gbm,xgboost}, \texttt{HistGradientBoostingClassifier} from scikit-learn
\cite{sklearn}; the linear baseline is \texttt{LogisticRegression} after a
\texttt{ColumnTransformer} one-hot encoding the 15 features. Hyperparameters are given in
the released code. Histogram binning \cite{lightgbm} keeps refitting across six substances, two
strategies and hundreds of bootstrap resamples affordable on CPU, and handles missing values
natively. \texttt{XGBClassifier} \cite{xgboost} at matched settings on the two largest
substances agrees to within 0.002 AUROC, so no conclusion depends on the boosting library.

All 15 TEDS fields are integer-coded categoricals. GBT receives these codes directly, since binning is
invariant to nominal-code ordering, whereas logistic regression needs the one-hot expansion. No scaling, imputation, or class rebalancing is applied.
We compare two training strategies that differ only in which episodes enter training:
\emph{substance-specific}, fitted and tested within a single substance, and \emph{pooled},
fitted once on all six high-volume substances and applied everywhere.

\textbf{Training assembly.} Every cross-sectional result comes from one fitted model. The
6{,}501{,}243 big-6 episodes are split 70/30 stratified jointly on primary substance and outcome,
giving 4{,}550{,}870 training and 1{,}950{,}373 held-out episodes, and the training side is then
reduced to 999{,}999 by subsampling on the same joint key. A single GBT is fitted on that
cohort and scored, without refitting, both on the big-6 holdout and on all 215{,}983 unseen-class
episodes. The temporal arms are year-restricted refits of the identical specification, training
filtered to 2015--2018 for the COVID arm and to years before 2022 for the temporal-coverage arm.
No unseen-class episode enters any training set, which is what makes the twelve held-out classes
out-of-sample by construction.

We use \emph{zero-shot} in the literal sense that the model sees zero training episodes of
the target substance class; the target (dropout versus completion) is unchanged, so the shift
at test time is in the covariate distribution rather than the label space. Because substance
is never a model input, the fitted model applies to a held-out class unmodified. All AUROCs carry bootstrap 95\% confidence intervals, and beyond discrimination we assess
calibration \cite{van_calster_calib} and clinical utility via decision-curve analysis
\cite{dca_vickers}. GBT and logistic experiments ran on one 48-core CPU machine; TabPFN-3
additionally used one NVIDIA RTX 2000 Ada GPU.


We also benchmark \textbf{TabPFN-3}, a pretrained tabular foundation model, under two
protocols. It predicts by in-context learning rather than parameter fitting, so cost scales with
(train~$\times$~test) size and the usable context is bounded by compute rather than data
availability: $\sim$10{,}000 rows on CPU versus 100{,}000 rows in 38.9~s on one GPU, a $60\times$ speedup at
matched size. Under \emph{matched training size} both
models receive the same 10{,}000 training and 2{,}000 test rows, isolating sample efficiency;
under \emph{maximum training size} each receives as much as it can use (100{,}000 rows for
TabPFN-3, up to $10^6$ for GBT), which asks which model to deploy. Both protocols score
identical test rows per comparison, so differences are paired, with 95\% bootstrap intervals
from 300 resamples.

\section{Results}

\subsection{Single-model coverage of unseen substances (primary result)}
Figure~\ref{fig:coverage} shows the headline result. The pooled model, trained only on the six
big-6 substances, predicts dropout for each of the twelve held-out small classes at a mean AUROC
of \textbf{0.777} (unweighted; range 0.745--0.806), so every small class is predictable despite
contributing nothing to training. Performance is not uniform: six classes sit above the big-6 internal reference (0.777) and six
below, the lowest being benzodiazepines (0.745), barbiturates (0.748) and OTC (0.751). No
class collapses toward chance. A temporal version of the
test (train on years $<$2022, evaluate on 2022 small-class episodes) gives a similar mean AUROC of
0.785, so coverage is not an artifact of same-period evaluation.

\begin{figure}[t]
\centering
\includegraphics[width=\columnwidth]{fig1_coverage.png}
\caption{\textbf{Coverage.} A pooled model trained only on the six big-6 substances predicts
dropout for twelve held-out small classes it never saw (mean AUROC 0.777; dashed line = big-6
internal reference). Bars are bootstrap 95\% CIs.}
\label{fig:coverage}
\end{figure}

\subsection{Pooling costs no accuracy}
For GBT, mean AUROC across the six highest-volume substances is 0.785 with a model trained on
the target substance alone and 0.776 with the pooled model, a difference of $-0.008$. The penalty is
small and consistent rather than concentrated in one substance: pooled training loses between
0.004 (Alcohol) and 0.012 (Other opioids) AUROC. As a supplementary analysis, rebuilding the pooled
cohort with an equal number of episodes per substance rather than matching the population mixture
moves per-substance AUROC by at most 0.004 (mean $+0.001$; it helps the four smaller substances
and costs the two largest), so the alcohol-heavy mixture is not the source of the pooling cost.
Pooling therefore does \emph{not} beat
substance-specific training, and the pooled model's value is coverage and single-model
deployment simplicity rather than an accuracy gain. Logistic regression reproduces the same
ordering at a uniformly lower level (0.753 and 0.742), trailing GBT by 0.031 to 0.034 AUROC,
which indicates nonlinear feature interactions matter.

\subsection{Calibration and clinical utility on unseen classes}
A risk score is actionable only if its probabilities mean what they say and if acting on them
beats the default policy. On the big-6 holdout, our reference, the pooled GBT reaches AUROC 0.791
and Brier 0.179.

Scoring the same model on all 215{,}983 unseen-class episodes gives AUROC 0.778 and Brier 0.186,
and calibration transfers essentially intact (Fig.~\ref{fig:calib}a): across risk
deciles the mean absolute gap between predicted and observed dropout is 0.015 and never exceeds
0.046, despite the unseen cohort's higher base rate (0.404 versus 0.397) and different substance
mix. The model is therefore usable as a probability, not merely a ranking, on classes it never saw. Decision-curve analysis (Fig.~\ref{fig:calib}b)
shows net benefit above both defaults through threshold 0.80: 0.275 at 0.20 and 0.126 at 0.50,
where ``treat all'' has already turned negative (crossing zero at 0.41). The margin over ``treat all'' widens above 0.15, the
regime that matters for a program supporting part of its caseload.

\begin{figure}[t]
\centering
\includegraphics[width=\columnwidth]{fig3b_calib_dca_unseen.png}
\caption{\textbf{Calibration and clinical utility on the twelve unseen classes}
($n=215{,}983$; pooled GBT trained only on the big-6). (a) Predicted versus observed dropout by
risk decile: the ten points are deciles of predicted risk pooled across all twelve classes
(about 21,600 episodes each), not individual substances, and the dashed diagonal is the
identity line on which predicted equals observed. Mean absolute deviation 0.015.
(b) Decision curve: net benefit stays above both
defaults through threshold 0.80, while ``treat all'' turns negative at 0.41.}
\label{fig:calib}
\end{figure}

\subsection{COVID-aware temporal transportability}
A model deployed on emerging substances also runs in a later era than the one it trained in, so
temporal and substance shift compound. The dropout base rate rises from $\sim$0.39 pre-2020 to
0.427 at COVID onset and stays elevated (0.426 in 2022), making the pandemic a natural shift test.
Training on 2015--2018 only, we hold 2019 out and evaluate on 2019, 2020 and 2022
(Fig.~\ref{fig:covid}).

On the trained substances the model is close to era-invariant: AUROC 0.786, 0.769, and 0.768:
a one-time step of 0.017 at the pandemic boundary and no further erosion. Repeating the design
end-to-end within the unseen classes (training on 125{,}912 unseen-class episodes from 2015--2018,
testing on 35{,}424 / 26{,}861 / 27{,}786) gives 0.789, 0.738, and 0.708: a 0.051 step at onset
and a further 0.030 by 2022, against 0.017 and 0.001. Combined substance and temporal shift
therefore costs about 0.081 AUROC over seven years, whereas either alone costs under 0.020.
Coverage survives the pandemic break but decays with time in a way same-period evaluation would
not reveal, arguing for periodic recalibration rather than indefinite deployment.

\begin{figure}[t]
\centering
\includegraphics[width=\columnwidth]{fig4b_covid_unseen.png}
\caption{\textbf{Temporal transportability across the COVID discontinuity.} Both models are
trained on 2015--2018 only, 2019 held out cleanly. (a) Test-year AUROC; error bars are
bootstrap 95\% CIs (big-6 only, where per-year test sets exceed 700k episodes). (b) AUROC lost
between consecutive eras: the unseen classes lose more at onset (0.051 versus 0.017) and keep
decaying afterwards (0.030 versus 0.001).}
\label{fig:covid}
\end{figure}

\subsection{Subgroup fairness}
Because algorithmic risk scores can encode and amplify demographic bias
\cite{obermeyer_bias}, we audit the model across subgroups. On the trained substances the audit is
clean: calibration gaps stay within $\pm$0.022 everywhere and age shows a monotonic gradient
(0.728 at ages 12--14 to 0.837 at 65+).

The audit that matters for our claim is on the unseen classes, across 22 race, sex and age strata
(Fig.~\ref{fig:fairness}). AUROC spans 0.677--0.798 (standard deviation 0.032), sex
is near-equal (male 0.775, female 0.780), age climbs from 0.707 at 12--14 to 0.795 at 55--64, and
the lowest values fall on the smallest strata (Alaska Native, $n{=}260$).

Calibration on unseen classes is where the audit finds a real exception. Twenty-one of 22 strata
stay within 0.045 of their observed rate, but the 65+ band predicts risk above observed by
\textbf{0.140} ($n{=}5{,}061$, observed dropout 0.165 against
a cohort rate of 0.404), so a threshold-based referral rule would over-flag this band. The gap is
localised rather than an age effect: a single state's other-tranquilizer episodes at 65+
($n{=}3{,}234$, 1.5\% of the cohort) report observed dropout 0.067 against 0.277 for big-6 patients
of the same age in the same state, a pattern absent from the training data; excluding them the 65+
band instead under-predicts slightly, by 0.010, and no stratum exceeds 0.05. Across the eight strata with $n>20{,}000$, calibration error never exceeds 0.017 and AUROC stays
in 0.772--0.798, so the model is fair in discrimination and adequately calibrated for the bulk of
the population.

\begin{figure}[t]
\centering
\includegraphics[width=\columnwidth]{fig5b_fairness_unseen.png}
\caption{\textbf{Subgroup audit on the twelve unseen classes} (22 strata; dashed line:
unseen-cohort AUROC 0.778). (a) AUROC by age band, race/ethnicity, and sex; the lowest values fall
on the smallest strata. (b) Absolute calibration gap per stratum: only the 65+ band (highlighted)
exceeds 0.05, at 0.140, while 21 of 22 stay within 0.045.}
\label{fig:fairness}
\end{figure}

\subsection{Tabular foundation model benchmark}
The two protocols give different answers, and the contrast is the result. At a matched
10{,}000/2{,}000-row train/test budget, TabPFN-3 beats GBT in \emph{every} one of the 18
substance-by-strategy comparisons by a mean of $+0.018$ AUROC. Given each model the data it can use,
that advantage disappears: TabPFN-3 at a 100{,}000-row context versus GBT at up to $10^6$ rows
wins only 8 of 18 comparisons, mean $-0.0001$ AUROC, with 15 of 18 paired bootstrap intervals
including zero, per-strategy means flat, and the three significant intervals pointing both ways.

TabPFN-3's advantage is therefore \emph{sample efficiency}, not higher attainable accuracy.
Both architectures converge near 0.78 AUROC, locating the limit in the intake-time
administrative variables rather than the learner, consistent with pooling also moving accuracy
very little.

\section{Discussion}
The central finding is that a single model covers substance classes it never trained on, which is
useful precisely where per-substance modeling fails: the long tail of substances too small
to support their own model. The dominant predictors (treatment setting, diagnosis
group, prior-treatment history) appear substance-general enough to serve all six,
supported by the flat pooled-versus-specific comparison and the coverage result.

\textbf{Limitations and future work.} The unit is a treatment episode, not a person, and the
outcome is an administrative discharge code, not a clinical assessment of recovery; 2021 is
excluded. The legitimate use is population-level (ranking a caseload to target finite retention resources),
while person-level deployment would require prospective, linked, IRB/DUA-governed validation.
Coverage here is zero-shot generalization, the weakest form of transfer; since pooling does not
beat substance-specific training we do not expect stronger adaptation to help, though a
confirmatory control (importance-weighted boosting \cite{importance_weight} or TabPFN few-shot
adaptation) is future work. Two findings bound the deployment claim. The 65+
calibration gap on unseen classes (0.140) is the one audit failure, and it traces to one
state-substance reporting cell rather than to age, so the indicated check is per-state label
screening, not recalibration;
temporal decay on unseen classes is larger than on trained substances and does not plateau, so the
result licenses periodic refitting rather than indefinite use of a frozen model.

\section{Conclusion}
On 6.72M national treatment episodes, a single pooled model predicts dropout for twelve
substance classes it never saw at mean AUROC 0.777. Testing deployment assurances \emph{on those
unseen classes} rather than on the training substances, three of four hold: calibration transfers
(mean decile gap 0.015), net benefit stays positive to threshold 0.80, and discrimination is
stable across strata (0.677--0.798); temporal robustness is qualified, since the COVID boundary
costs 0.051 AUROC on unseen classes against 0.017 on trained ones. Subject to periodic
recalibration and an age-specific check, one model can cover the substance long tail, including
substances that emerge after training. TabPFN-3 is the more sample-efficient learner, reaching
the same accuracy band from a tenth of the episodes, but at each model's maximum training size the
two are indistinguishable, locating the limit in the intake data, not the learner.

\section*{Reproducibility and Data Availability}
TEDS-D files are public from SAMHSA; all models use scikit-learn implementations and the public
TabPFN-3 checkpoint. The code repository contains the analysis scripts, hyperparameters, a data
dictionary and one CSV of results per experiment, including every value plotted here.

\begin{thebibliography}{00}
\bibitem{chekroud_ml_psych} A. M. Chekroud \textit{et al.}, ``The promise of machine learning in predicting treatment outcomes in psychiatry,'' \textit{World Psychiatry}, vol.~20, pp.~154--170, 2021.
\bibitem{ml_teds_disparities} A. Baird, Y. Cheng, and Y. Xia, ``Use of machine learning to examine disparities in completion of substance use disorder treatment,'' \textit{PLoS ONE}, vol.~17, p.~e0275054, 2022.
\bibitem{ml_jamia} M. Nasir \textit{et al.}, ``Machine learning-based outcome prediction and novel hypotheses generation for substance use disorder treatment,'' \textit{J.\ Amer.\ Med.\ Inform.\ Assoc.}, vol.~28, pp.~1216--1224, 2021.
\bibitem{attrition2025} J. A. Rabinowitz \textit{et al.}, ``Predictors of treatment attrition among individuals in substance use disorder treatment: A machine learning approach,'' \textit{Addict.\ Behav.}, vol.~163, p.~108265, 2025.
\bibitem{sud_ml_review} B. P. de Mattos \textit{et al.}, ``Craving for a robust methodology: A systematic review of machine learning algorithms on substance-use disorders treatment outcomes,'' \textit{Int.\ J.\ Ment.\ Health Addict.}, vol.~24, no.~2, pp.~1090--1117, 2026, doi: 10.1007/s11469-024-01403-z.
\bibitem{tabpfn_nature} N. Hollmann \textit{et al.}, ``Accurate predictions on small data with a tabular foundation model,'' \textit{Nature}, vol.~637, pp.~319--326, 2025.
\bibitem{bailey2022} J. D. Bailey and A. DeFulio, ``Predicting substance use treatment failure with transfer learning,'' \textit{Subst.\ Use Misuse}, vol.~57, pp.~1982--1987, 2022.
\bibitem{friedman_gbm} J. H. Friedman, ``Greedy function approximation: A gradient boosting machine,'' \textit{Ann.\ Statist.}, vol.~29, pp.~1189--1232, 2001.
\bibitem{xgboost} T. Chen and C. Guestrin, ``XGBoost: A scalable tree boosting system,'' in \textit{Proc.\ ACM SIGKDD}, 2016, pp.~785--794.
\bibitem{sklearn} F. Pedregosa \textit{et al.}, ``Scikit-learn: Machine learning in Python,'' \textit{J.\ Mach.\ Learn.\ Res.}, vol.~12, pp.~2825--2830, 2011.
\bibitem{lightgbm} G. Ke \textit{et al.}, ``LightGBM: A highly efficient gradient boosting decision tree,'' in \textit{Adv.\ Neural Inf.\ Process.\ Syst.\ (NIPS)}, 2017, pp.~3146--3154.
\bibitem{van_calster_calib} B. Van Calster \textit{et al.}, ``Calibration: The Achilles heel of predictive analytics,'' \textit{BMC Med.}, vol.~17, p.~230, 2019.
\bibitem{dca_vickers} A. J. Vickers and E. B. Elkin, ``Decision curve analysis: A novel method for evaluating prediction models,'' \textit{Med.\ Decis.\ Making}, vol.~26, pp.~565--574, 2006.
\bibitem{obermeyer_bias} Z. Obermeyer \textit{et al.}, ``Dissecting racial bias in an algorithm used to manage the health of populations,'' \textit{Science}, vol.~366, pp.~447--453, 2019.
\bibitem{importance_weight} N. Lu \textit{et al.}, ``Rethinking importance weighting for transfer learning,'' arXiv:2112.10157, 2021.
\end{thebibliography}
\end{document}

